% Canonical node 014 · C014; current section path 1.2.7.
% Canonical owner: C014
\cnnaNodeHeading{014}{C014}{Bootstrap state $X_1$}
\cnnaNodePlacement{1.2.7}{D1}

\paragraph*{Position and scientific role.}
\texttt{C014} closes the exceptional bootstrap section.  It packages the
concrete C013 birth and imports the two independent certificates T001 and M005.
The result is the base case handed to the recurrent state schema \texttt{C005}.

% CNNA-ARCHITECTURE-BEGIN C014
\begin{cnnaInfoBox}{First response-capable CNNA state}
C014 packages the first weighted two-node net.  From this point onward the complement architecture can become dynamical: later nodes may retain part of the net, eliminate a complementary interior, read the resulting response, and feed it into the next birth.  C014 itself performs none of those later operations; it supplies their minimal nontrivial carrier.
\end{cnnaInfoBox}
% CNNA-ARCHITECTURE-END C014

\paragraph*{Explicit construction.}
The mathematical contents are
\begin{equation}
  X_1=\bigl(V_1,\mathcal R_1,\mathcal C_1\bigr),
  \qquad
  V_1=\{r,v_1\},
  \quad
  \mathcal R_1=\bigl((r,v_1),(v_1,r)\bigr),
  \quad
  \mathcal C_1=(1,1),
  \label{eq:c014-bootstrap-state}
\end{equation}
with $r=\varepsilon$ and $v_1=(0)$.  In the implementation this tuple is not
stored redundantly: \texttt{BootstrapState} contains exactly one field, the
verified C013 birth, and exposes the endpoints, relations and conductances by
projection.

\paragraph*{Invariants and certificates.}
C014 inherits the root, newborn and unit-pair invariants from C013.  Lean then
lifts T001 through the wrapper, proving that changing the explicit seed leaves
$X_1$ unchanged.  It also rationally lifts the two actually stored conductances,
identifies both with the N001 representative, and applies M005 separately to
each orientation.  The comparison scale is not stored in $X_1$ and no second
bootstrap state is constructed.

\paragraph*{Boundary cases and counterchecks.}
C014 can be built only after C013, hence only when $L\ge1$.  At $L=0$ the
root-only approximant exists but no $X_1$ does.  Python rejects a non-C013
payload and verifies that the dataclass has only the field \texttt{birth}.
Neither layer stores the seed, a variable unit, a response value, a time,
geometry or recurrent-update data.  ``Response-capable'' means only that a
nontrivial weighted relation now exists on which later response operators can
act.

\paragraph*{Cross-layer agreement, result and handoff.}
Python realizes the minimal wrapper and executable projections.  Lean realizes
the same wrapper and additionally attaches theorem-level seed and unit
certificates without adding runtime fields.  Verified edge \texttt{E014}
provides the recurrent base case to \texttt{C005}.  The later finite-iteration
handoffs \texttt{E077} and \texttt{E153} remain planned and blocked; C014 by
itself proves no recurrent existence theorem.

\noindent\textbf{Primary code anchors.}
Python state, projections, constructor and test: lines 23--51 and 18--30.
Lean wrapper, projections, seed-neutrality lift, rational conductance lift and
unit-representative theorem: lines 23--99.  The Supplement lists all 12 anchors
with exact source hashes.
